\documentclass[12pt,a4paper]{article}
\usepackage[utf8]{inputenc}
\usepackage{geometry}
\geometry{left=2.5cm,right=2.5cm,top=2.5cm,bottom=2.5cm}
\usepackage{hyperref}

\title{\textit{Being is Time} (Yixin Jing): An Introduction}
\author{Zhang Yu}
\date{Chinese edition published 2024; English edition forthcoming}

\begin{document}
\maketitle

\section*{Overview}
\textit{Being is Time} (Chinese title: \textit{Yixin Jing}) is a philosophical and scientific treatise that unifies physics, economics, and metaphysics under a single ontological framework: the 2D tension surface. The book consists of nine parts, encompassing 27 chapters, and lays the foundation for the N.E.A. (Network Emergent Architecture) Framework. The Chinese edition (29 chapters) was published in 2024; an English edition is forthcoming in 2026.

\section*{Core Philosophical Thesis}
The central thesis is \textbf{Being is Time}. Space, matter, and forces are not fundamental; they are emergent projections of a 2D information surface (the “tapestry”) whose tension and refresh rate generate the perceived three‑dimensional world. Time is the logical refresh rate of this surface; gravity is the gradient of its tension; quantum mechanics is the discrete exchange of its constituent bits.

\section*{Structure of the Book}
The book is organized into nine parts:

\begin{itemize}
    \item \textbf{Part I: Transaction} – Introduces entropy, value, price, and the triple model (Chapters 1–5).
    \item \textbf{Part II: Scholarship} – Analyzes entropy across mathematics, physics, economics, and biology (Chapters 6–10).
    \item \textbf{Part III: People} – Applies the framework to history, religion, politics, and economics (Chapters 11–16).
    \item \textbf{Part IV: Forward} – Discusses future stable states and technological evolution (Chapters 17–18).
    \item \textbf{Part V: Retrospect} – Explores physics, philosophy, and theology (Chapters 19–20).
    \item \textbf{Part VI: Abyss} – Builds a layered cosmos (space, force, life, wisdom) and introduces the concepts of “tension surface” and “dimensional folding” (Chapters 21–22).
    \item \textbf{Part VII: Light and Shadow} – Reinterprets time, space, and the universe through a photon‑based cosmology (Chapters 23–24).
    \item \textbf{Part VIII: The One} – Concludes with dimensional unification and the nature of reality (Chapters 25–26).
    \item \textbf{Part IX: Nine Nine} – A final synthesis, unifying all concepts (Chapter 27).
\end{itemize}

\subsection*{Chapters 22–29 and Their Relevance to the N.E.A. Papers}
Chapters 22 through 29 contain the most advanced physical and metaphysical speculations of the book, many of which directly inspired the N.E.A. research program:

\begin{itemize}
    \item \textbf{Chapter 22 (The Perfect Tapestry)} introduces the “2D tension surface” as the fundamental substrate, coins the term “Yi-Xin” (易心) and “Zhang-Suo Sheng-Ceng” (张缩生层 — expansion/contraction layering). It proposes that space is a Bose–Einstein condensate of “space quanta” and that time emerges from the sequential interaction of photons with these quanta. This chapter lays the philosophical groundwork for the holographic and emergent gravity ideas used in the N.E.A. papers.
    \item \textbf{Chapter 23 (Light)} explores the nature of light, the Rindler horizon, and the “Babel tower” metaphor, which later influenced the numerical Rindler experiments in Paper IV.
    \item \textbf{Chapter 24 (Shadow)} presents a multi‑layered universe model (Space Sea, Love Sea, Force Sea, Life Sea, Wisdom Sea) and the idea that gravity is a collective entropic force. This directly anticipates the entropic gravity interpretation adopted in the N.E.A. framework.
    \item \textbf{Chapter 25 (All Methods Return to One)} discusses dimensions, numbers, and idealism, ending with the conclusion that expansion and contraction are the only fundamental operations.
    \item \textbf{Chapter 26 (Falling in Love with Your Form)} is a long, dense chapter that re‑derives all fundamental physics from the “Yi-Xin” principles: it derives momentum conservation from spatial translation symmetry, energy conservation from time translation symmetry, and identifies the Lagrangian and Hamiltonian with “Yi” (change) and “Xin” (invariant). It then reconstructs mass as “entangled spacetime”, weak and electromagnetic forces as folded spacetime, and the strong force as the hidden third fold. This chapter is the direct predecessor of the unified ontology (Section 6) in the N.E.A. gravity paper.
    \item \textbf{Chapter 27 (Being is Time)} provides the final philosophical closure, linking the mathematical structure (Noether’s theorem, least action principle) to the metaphysical insight that all existence is a manifestation of time. It also contains the author’s personal reflections and a discussion of the limits of mathematics.
\end{itemize}

\section*{Relation to the N.E.A. Papers}
The N.E.A. (Network Emergent Architecture) research program comprises four papers (Papers I–IV) that formalize, test, and extend the ideas first presented in \textit{Being is Time}. 

\begin{itemize}
    \item \textbf{Paper I} (\textit{From Rigid Axioms to Emergent Enthalpy}) takes the philosophical “Yi-Xin” duality and casts it into two rigid axioms (Noether’s theorem and the principle of least action). It shows that the enthalpy decomposition \(Q = U + pV\) (physics) and \(Q = k N \ln I + pE(N)\) (economics) are emergent consequences of these axioms under specific computational constraints—a formalization of the “expansion/contraction layering” concept.
    \item \textbf{Paper II} (\textit{The Physical Universe from the N.E.A. Framework}) builds on the physical enthalpy to derive the uniqueness of three dimensions, the Einstein field equations as a thermodynamic identity, and the gauge group \(U(1) \times SU(2) \times SU(3)\) from spectral degeneracy. This paper gives a rigorous, graph‑theoretic foundation to the “dimensional folding” and “tension surface” speculations of Chapters 22 and 26.
    \item \textbf{Paper III} (\textit{Economic Order and Social Cycles}) applies the same framework to economics, deriving the Pareto distribution and the cycles of inequality—a quantitative version of the “cycles of order and chaos” described in Parts III and IV of the book.
    \item \textbf{Paper IV} (\textit{From Holographic Dimension to Tension-Induced Gravity}) presents the observational evidence: the N.E.A. model explains galaxy rotation curves without dark matter, with the effective dimension \(q\) transitioning from 2 (3D Newtonian) to 1 (2D gravity) in low‑density regions. This paper also includes a unifying ontology (Section 6) that directly synthesizes the philosophical insights of Chapters 22–27 into a coherent physical picture—special relativity as kinematics on the 2D surface, general relativity as tension dynamics, and quantum mechanics as discrete bit‑exchange.
\end{itemize}

Thus, the book provides the philosophical and conceptual bedrock; the N.E.A. papers provide the mathematical formulation, numerical experiments, and empirical tests that validate and extend those ideas. Together they form a unified body of work that reinterprets physics, economics, and cosmology from a single informational perspective.

\section*{Key Contributions of the Combined Work}
\begin{itemize}
    \item \textbf{Unification of special relativity, general relativity, and quantum mechanics} within a single 2D substrate (tension surface).
    \item \textbf{Dark matter as a projection artifact} of dimensional collapse (\(q \to 1\)).
    \item \textbf{Explanation of the cosmological constant} via area‑law vacuum energy.
    \item \textbf{Fermion spectrum as topological defects} of the 2D surface under 4D anchoring.
    \item \textbf{Economic cycles and inequality} as maximum‑entropy outcomes under finite‑bandwidth constraints.
\end{itemize}

\section*{Citation}
If you refer to the book in academic work, please cite as:
\begin{verbatim}
Zhang, Y. (2024). Being is Time (Yixin Jing). [Chinese edition]. 
English edition forthcoming, 2026.
\end{verbatim}

For the N.E.A. papers, please cite the corresponding preprint series (see the paper’s bibliography).

\section*{Availability}
The Chinese edition of \textit{Being is Time} is available in print and digital formats. The English edition will be released in 2026. For updates, visit the author’s website or the GitHub repository of the N.E.A. project.

\end{document}